\section{Introduction: The Oldest Dream of Domestication}
\label{sec:intro}

For at least fifteen thousand years---and by some accounts more than
thirty thousand---humans have lived alongside dogs, and for nearly as long we
have wanted to know what they are saying. The wish is older than science:
folktales grant the hero a magic fern-seed or a drop of dragon's blood that
lets him understand the speech of beasts. What is new is that the wish now
arrives packaged as a product. A smartphone application records a bark and
returns a sentence---``I want to play,'' ``my tummy hurts,'' ``I love
you''---in under a second. The vendor advertises an accuracy. The number is
large. And the curious reader is left with a question that turns out to be
surprisingly hard to answer: \emph{how far are we, actually, from translating
what a dog says?}

This survey takes that question seriously and refuses to answer it with a
single number. ``Translation'' is not one capability but a stack of them, and
the distance to each layer is different. A system that reports whether a dog
is relaxed or distressed is doing something real and useful; a system that
claims to render the dog's ``words'' as an English sentence is doing
something else, and the difference between the two is not a matter of degree
but of kind. Our central move is to make that stack explicit---a five-rung
\emph{translation ladder} (Section~\ref{sec:ladder})---and to locate current
research and current products on it. The headline finding is easy to state
and, we will argue, well supported: \textbf{today's best systems live on the
second rung, affect classification, with a handful of careful laboratory
studies reaching toward the third.} Everything above that remains either
aspirational or, in the strong sense, ill-posed.

\subsection{Why now: three converging currents}
Three developments make this the right moment for such a survey. First,
\emph{bioacoustic foundation models} have matured. Self-supervised encoders
trained on animal sound---AVES~\cite{hagiwara2023aves}, the audio--language
model NatureLM-audio~\cite{robinson2024naturelm}, and Google DeepMind's Perch
line~\cite{ghani2023perch}---now transfer across species and tasks, and
standardized benchmarks such as BEANS~\cite{hagiwara2023beans} and BirdSet
give the field common measuring sticks. Second, \emph{flagship decoding
projects} have produced concrete, peer-reviewed structure discoveries:
Project CETI's sperm-whale ``phonetic alphabet''~\cite{sharma2024whale},
Google's DolphinGemma sequence model~\cite{google2025dolphingemma}, and the
Earth Species Project's interactive songbird work. Third, \emph{canine
cognition} has its own empirical renaissance, from the ELTE Budapest group's
decades of bark and growl psychophysics~\cite{pongracz2005barks,farago2017growls}
to the UC San Diego soundboard studies that put ``talking buttons'' under
controlled observation~\cite{bastos2024plos,bastos2024scirep}. The
intersection of these three currents is exactly the space this paper maps.

\subsection{The thesis in one paragraph}
A dog's communicative repertoire is rich, multimodal, and genuinely
informative---but it is, in the technical sense, a \emph{signaling system}
rather than a \emph{language}. It encodes affect, arousal, motivational
state, individual identity, and some referential content (a specific toy, a
specific door), yet it lacks the open-ended, recursive, generative syntax
that human language translation presupposes. Consequently, ``translating
what a dog says'' is, above the level of affect and intent, partly a
\emph{category error}: there may be no sentence-shaped meaning in the signal
to recover. This is not a counsel of despair. It reframes the achievable
goal---validated, calibrated affect-and-intent inference---and exposes why
the field's most quotable numbers (``81.5\% accuracy,'' ``94.6\% accuracy'')
are, as presently defined, close to meaningless: there is no ground-truth
intended message to score a translation against. We call this the
\emph{falsifiability problem}, and it is the methodological heart of the
survey (Section~\ref{sec:eval}).

\subsection{Contributions and organization}
This paper makes four contributions.
\begin{enumerate}[label=(\arabic*)]
  \item \textbf{A conceptual ladder} that separates signal, affect, intent,
  reference, and dialogue, and uses it to adjudicate over-claiming
  (Section~\ref{sec:ladder}, Figure~\ref{fig:ladder}).
  \item \textbf{A unified technical synthesis} of the decoding pipeline---%
  encoding, structure discovery, grounding, validation---connecting canine
  work to the broader cross-species literature
  (Sections~\ref{sec:biology}--\ref{sec:crossspecies}).
  \item \textbf{A critical audit} of commercial ``pet translators'' and the
  evaluation crisis that surrounds them
  (Sections~\ref{sec:products}--\ref{sec:eval}).
  \item \textbf{An honest distance estimate and roadmap}, with confidence
  bands and an ethics framework
  (Sections~\ref{sec:distance}--\ref{sec:future}, Figure~\ref{fig:roadmap}).
\end{enumerate}
Section~\ref{sec:biology} grounds the discussion in what dogs actually
signal; Section~\ref{sec:conclusion} concludes. Throughout, we keep two
registers separate: what is \emph{measured} (peer-reviewed effect sizes,
model metrics on defined tasks) and what is \emph{claimed} (vendor
accuracies, press-release superlatives). The gap between those registers is,
in a real sense, the subject of this paper.